\section{Wavelength calibration}
\label{sec:wavelength}

\subsection{How \banzai{} calibrates wavelengths}
\label{sec:wl_banzai}
For each arc, \banzai{} collapses each order along an assumed line tilt, finds and matches peaks, centroids bright
isolated lines row by row with a Gauss--Hermite profile \citep{Cappellari2017}, fits blends with a fixed line-spread
function, fits a straight line $x=a+b\,y_{\rm o}$ per line (the tilt), and finally fits
$x(\lambda)$ as a Legendre series (degree 5 in the red order, 2 in the blue in the current code; the 1.0.1 products
carry degree 4 and 2) and inverts it to $\lambda(x)$. The tilt is a constant ($\approx8.2^\circ$) at COJ. There is no
outlier rejection in the final fit, and no per-exposure correction of the zero point (flexure or position in the
slit): every science frame inherits the solution of its arc.

\subsection{The arc line list leaves the RM lines unconstrained}
\label{sec:wl_lines}
Of the 68 catalogue lines, 33 are used. In the blue order the solution rests on Hg lines only: 3650/3663, 4047,
4078, 4358 and then 5461, 5770, 5791\,\AA{}; there is no line between 4358 and 5461\,\AA{}, exactly where H$\beta$ and
[O\,{\sc iii}] fall for $z\lesssim0.1$. The red order starts at $\approx4550$\,\AA{} but its first used line is
5461\,\AA{}, so 4550--5461\,\AA{} is extrapolated, and it has no line between 5791 and 6965\,\AA{} (H$\alpha$ for
$z\approx0.05$). The Zn lines that would fill the blue gap are flagged ``removed until we start turning on the Zn
lamp again''.

We collapsed all 64 raw arcs with the wavelength map of a science frame and aligned them by cross-correlation
(night-to-night shifts 1.07\,\AA{} rms red, 0.49\,\AA{} blue). Three facts follow (Fig.~\ref{fig:arcs}):
(i) Hg\,4916 is detected in every single 60\,s arc in both orders, and Ar\,6416, 6677 and 6753/6767 in the red, so
stacking is not needed; (ii) although the Zn lamp is ``on'' in 63 of 64 headers, no Zn line is detected even in the
stack of 64 arcs; (iii) with the 6\,arcsec slit all lines are $\approx51$\,\AA{} (red) and $\approx25$\,\AA{}
(blue) top-hats, so many neighbours blend.

\subsection{Arc refit with the unused lines}
\label{sec:wl_refit}
Each raw arc is overscan-subtracted, trimmed and collapsed with the frame's own 2-d wavelength map (which contains
the tilt) by averaging all pixels within 30 rows of the order centre in bins of half a pixel. Lines are measured in
\emph{groups} $g$ of laboratory wavelengths $\lambda_{g,j}$ (Table~\ref{tab:groups}) with the forward model
\begin{equation}
M_g(\lambda) = c_0 + c_1\frac{\lambda-\bar\lambda_g}{100} + \sum_j A_j\,B(\lambda;\lambda_{g,j}+s_g,W,\sigma),
\label{eq:arcgroup}
\end{equation}
\begin{equation}
B(\lambda;\mu,W,\sigma)=\frac{1}{2}\left[\mathrm{erf}\!\left(\frac{\lambda-\mu+W/2}{\sqrt2\,\sigma}\right)
 -\mathrm{erf}\!\left(\frac{\lambda-\mu-W/2}{\sqrt2\,\sigma}\right)\right],
\label{eq:boxcar}
\end{equation}
i.e. a top-hat of the slit-image width $W$ convolved with a Gaussian of width $\sigma$, one per line at the
laboratory separations, with a common shift $s_g$ (measured minus laboratory wavelength on the current solution)
and free amplitudes $A_j\ge0$. The initial width is $W=8.5$\,\AA\,arcsec$^{-1}$ (red) or
4.25\,\AA\,arcsec$^{-1}$ (blue) times the slit width. Pairs closer than $0.3W$ whose catalogue strengths are known
are merged into one line at their strength-weighted mean wavelength (their amplitude ratio is otherwise degenerate
with $s_g$); if a strength is unknown, the components are kept with free amplitudes. A group is accepted if the fit
converged, $|s_g|<14.5$\,\AA{} and its peak exceeds 10 times the residual noise.

The correction is a weighted Legendre series fitted to the group shifts,
\begin{equation}
\Delta\lambda(\lambda)=\sum_{k=0}^{K}d_k\,L_k(\tilde\lambda),\qquad
w_g = \left(\sigma_{s_g}^2+0.3^2+\delta_g^2\right)^{-1},
\label{eq:arccorr}
\end{equation}
with $K=5$ (red) and 2 (blue), $\tilde\lambda$ the wavelength mapped to $[-1,1]$, $\sigma_{s_g}$ the statistical error
from the Jacobian, 0.3\,\AA{} a floor and $\delta_g$ an extra systematic for unresolved blends (0.5--1\,\AA{},
Table~\ref{tab:groups}); groups deviating by more than $4\sigma$ are iteratively rejected. Outside the span of
the accepted groups the correction is continued linearly from its end values and slopes, so that a degree-5 fit never
extrapolates wildly. The corrected wavelength is $\lambda'=\lambda-\Delta\lambda(\lambda)$.

The per-night fits have residuals of 0.48\,\AA{} rms (red, median of 23 groups) and 0.10\,\AA{} (blue, 4 groups).
The red-order error is $\approx-8$\,\AA{} at Hg\,4916 and $\approx-3$\,\AA{} at 5240\,\AA{}
(Fig.~\ref{fig:arcs}, bottom). The arc solution is stable in shape: after removing each night's zero point the
solutions of the 50 frames agree to $<0.4$\,\AA{} over most of each order. We therefore apply the \emph{median} of the
nightly corrections to every frame; using each night's own fit adds fit noise (Section~\ref{sec:results_wl}).

\subsection{The zero point of the object's light}
\label{sec:wl_zp}
The arc and the night sky fill the slit; a compact object does not. In a wide slit the object's wavelength zero
point depends on where it sits across the slit: a displacement $\Delta\theta$ moves its spectrum by
$\Delta x=\Delta\theta/0.337''$ pixels, i.e. $\approx10$\,\AA{} per arcsec in the red order. Neither the arc nor the
sky lines can measure this. The telluric bands can (Section~\ref{sec:diag_tell}): for each band $b$ we measure
$s_b$ (eq.~\ref{eq:tell}), remove its offset $o_b$ (Table~\ref{tab:bands}; calibrated on the 18 EG\,21 spectra after the
arc refit), and combine
\begin{equation}
s_{\rm red}=\frac{\sum_b \omega_b\,(s_b-o_b)/(\sigma_{s_b}^2+0.5^2)}{\sum_b \omega_b/(\sigma_{s_b}^2+0.5^2)},
\label{eq:tellcomb}
\end{equation}
with weights $\omega_b$ (1 for A, 0.5 for B and H$_2$O\,7200, 0.25 for H$_2$O\,9300) and iterative rejection of
bands deviating by more than $3\sigma$. Both orders image the same slit, so the blue-order shift in pixels is the red
one times the ratio of the anamorphic magnifications:
\begin{equation}
s_{\rm blue} = \rho\,\frac{s_{\rm red}}{D_{\rm red}}\,D_{\rm blue},\qquad \rho=1.13,
\label{eq:transfer}
\end{equation}
with $D$ the dispersions (\AA\,pixel$^{-1}$). $\rho$ was measured once from the slope of the blue-order [O\,{\sc
iii}] pixel shift against the red-order A-band pixel shift (slope 1.13--1.14, $r=0.99$, residual 0.2\,pixel; it is
an instrument constant). The corrected wavelengths are $\lambda''=\lambda'-s_{\rm order}$. If no band is usable, the
sky zero point (Section~\ref{sec:diag_sky}) is used instead (this corrects flexure but not the slit position).

The absolute zero point is that of the telluric template. It is uncertain by $\approx\pm4$\,\AA{}: on the
standards the A band lies $-3.5\pm0.7$\,\AA{} from the sky lines (if the standards are centred in the slit), whereas
the [O\,{\sc iii}] lines of Fairall~9 imply $+3.9$\,\AA{} (if [O\,{\sc iii}] is at the systemic redshift). This is a
single constant; it does not affect the consistency between exposures or between the two orders.
